\section{Related Work}

This section covers the literatures the article draws on, in the order the argument needs them. It begins with POMDP solvers and the $\rho$-POMDP framework that supplies the formal object of study, then turns to the classical value-of-information and experimental-design results on which the identification of reward with information rests. It then takes up active inference and its expected free energy objective, from which the information weight is derived, and the sophisticated-inference form the paper builds on. It closes with three neighboring lines, namely scaling active inference, exploration and intrinsic motivation, and resource-constrained information acquisition.

\paragraph{POMDPs and solvers.}
POMDPs formalize sequential decision-making under state uncertainty \citep{smallwood1973, kaelbling1998}. Exact finite-horizon POMDP planning is PSPACE-complete. Point-based offline methods (PBVI \citep{pineau2003}, HSVI \citep{smith2004}, SARSOP \citep{kurniawati2008}) approximate the value function on reachable beliefs \citep{shani2013}. Online solvers plan from the current belief. POMCP \citep{silver2010} uses Monte Carlo tree search (MCTS) with UCB1 action selection and evaluates leaves by rollouts. UCB1 is the upper-confidence-bound rule that picks the action whose estimated return plus an exploration bonus is largest. DESPOT \citep{ye2017} searches a regularized sparse belief tree, sparse because each node branches on a small set of sampled scenarios rather than on every possible observation, and regularized because the search penalizes the size of the policy it builds so as not to overfit those samples. POMCPOW \citep{sunberg2018} extends POMCP with progressive widening, a technique that limits how many children a search node expands as a function of its visit count, so that the tree can be grown over continuous spaces. All explore implicitly through stochastic simulations rather than explicitly valuing information gain. A recent counterpoint reasons about value of information explicitly inside the tree search. \citet{laouar2026voi}, sharing an author with \citet{sunberg2018}, propose Value of Information Monte Carlo planning, which selectively skips observation branches whose expected value of information is low rather than exploring them uniformly. Our EFE agent sits on the same side of that divide but values information in a different way. It computes the expected information gain of each observation action in closed form from the current belief, whereas POMCP estimates the value of a branch from sampled rollouts.

The rest of this paragraph previews where the two valuations are compared and what controls accompany the comparison. Section~\ref{sec:discussion} and Appendix~\ref{app:pomcp} report where the closed-form information valuation outperforms POMCP with semi-informed rollouts. A semi-informed rollout, defined in Appendix~\ref{app:pomcp}, picks observation actions at random but ends in a belief-optimal commit, the commit action that maximizes expected reward under the rollout's Bayesian-updated belief. The same section and appendix also test whether that advantage can instead be attributed to solver configuration, using three controls, and state what those controls do and do not isolate. The first is a per-environment sweep of POMCP's exploration constant, the coefficient on the UCB1 exploration bonus just described. The second is an information-gain rollout that, unlike the semi-informed one, picks the observation with the largest exact one-step information gain at every rollout step. The third is a component ablation of MCTS-EFE, the paper's own tree-search variant, which the paragraph on scaling active inference below introduces and that section defines. The ablation switches two of MCTS-EFE's components, the in-tree information-gain reward and the max-backup rule, on and off in every combination.

\paragraph{$\rho$-POMDPs.}
\citet{araya2010} introduced $\rho$-POMDPs, replacing the state-dependent reward $R(s,a)$ with a belief-dependent reward $\rho : \Delta(S) \times \mathcal{A} \to \mathbb{R}$, where $\Delta(S)$ denotes the set of probability distributions over states, so the objective becomes $\max_\pi \mathbb{E}_\pi[\sum_t \gamma^t \rho(b_t,a_t)]$. In their notation $\rho$ denotes the whole belief-dependent reward, and they note that it may itself combine the expected state-action reward with an uncertainty measure. This paper adopts that additive form but reserves the symbol $\rho$ for the belief-dependent information term alone, so that the per-step utility is written $R(s,a) + \rho(b,a)$, with $R$ the ordinary state-dependent reward. The information weight $w$ of the Introduction lives inside $\rho$. The $\rho$ this paper studies is $w$ times the expected information gain of an observation action. At $w{=}1$, the EFE case, it is exactly the $\rho_{\mathrm{EFE}}$ of Proposition~\ref{prop:equivalence}, the expected information gain of the chosen observation action and zero on commit actions. For general $w$ the budgeted sections write the same utility as $R + wI$, with the information gain abbreviated to $I$, since there the weight itself is the object of study.

Whether standard solver machinery carries over depends on the shape of $\rho$, and \citet{araya2010} showed when it does. When $\rho$ is convex in the belief, the value function is likewise convex. A convex $\rho$ can moreover be approximated arbitrarily well by a piecewise linear and convex (PWLC) function. Once $\rho$ is PWLC the value function itself is PWLC, so standard solvers can be reused with limited changes. \citet{fehr2018} extended this to Lipschitz-continuous non-convex $\rho$, and \citet{benchetrit2025} developed an anytime incremental $\rho$-POMDP planner for continuous spaces built on progressive widening. The $\rho$ studied in this paper is not among the convex cases. It does not inherit the convex value-function structure that convex $\rho$ yields \citep{araya2010}, nor the PWLC approximation that structure supports. Expected information gain is concave, not convex, in the belief, so $\rho_{\mathrm{EFE}}$, the symbol Section~\ref{sec:methodology} adopts for it, falls in the non-convex regime. For that regime \citet{fehr2018} give Lipschitz-continuous $\epsilon$-optimal value functions when $\rho$ is itself Lipschitz in the belief. Expected information gain meets that hypothesis whenever every observation likelihood is strictly positive, since $I_a(b) = H(O_a b) - \sum_s b(s)\, H(O_a(\cdot \mid s))$ with $O_a b$ the predictive observation distribution, which then stays a fixed distance from the boundary of the observation simplex where the entropy's slope is unbounded (for a binary distribution $H(\epsilon, 1-\epsilon)/\epsilon$ diverges as $\epsilon \to 0$). It can fail for observation models with zero entries, such as a noiseless sensor, where that predictive distribution reaches a vertex. None of this article's results relies on the guarantee either way. The exact tree enumeration of Section~\ref{sec:methodology} evaluates the recursion directly and needs no approximation guarantee.

The same framework-level move, rewarding information explicitly rather than only through its effect on state-dependent reward, also appears under another name. POMDPs with information rewards (POMDP-IR) \citep{spaan2015} augment planning with explicit information rewards. That framework was developed for active cooperative perception, in which sensing actions are chosen across cooperating sensors to serve a shared task. \citet{satsangi2018} exploit submodular value functions for scaling active perception and note the equivalence between POMDP-IR and $\rho$-POMDPs. The aim behind $\rho$-POMDPs predates their formalization. \citet{boutilier2002} casts preference elicitation as a POMDP whose hidden state is the user's utility function and whose belief is a distribution over those functions, selecting the next preference query by sequential rather than myopic value of information. Its reward remains a function of the state, so it is a predecessor in motivation rather than in construction. What it shares with $\rho$-POMDPs is the aim of choosing queries for what they reveal, with an information-gathering criterion fixed by that problem rather than derived from a variational objective.

Common choices for $\rho$ (entropy reduction, KL divergence from a target belief, and information gain) each encode a different notion of epistemic value, but the choice and weight remain heuristic and environment-specific. We instead derive the information-unit weight of the Introduction from the variational objective of active inference, yielding an untuned default under a shared reward convention. That convention, one reward unit per unit of information, is the exchange rate that fixes the coefficient. Observation costs are stated on that same reward scale. That unit of information can be taken in bits or in nats, and the paper uses both. The Introduction's information term is measured in bits, the base-two logarithm unit, while the EFE derivation is written in nats, the natural-logarithm unit, and the two differ by a constant factor. The paragraph on reward-to-nats calibration in Section~\ref{sec:methodology} makes this precise and states why the default is tied to that scale rather than invariant to it.

\paragraph{Value of information and experimental design.}
The idea that information has quantifiable decision-theoretic value predates both POMDPs and active inference. \citet{howard1966} formalized the value of information in decision analysis, and \citet{lindley1956} introduced expected information gain as a criterion for optimal Bayesian experimental design. One classical result grounds the identification of reward with information on which this paper rests. Under a proper log scoring rule, proper meaning that reporting one's true belief maximizes the expected score, the agent is paid the logarithm of the probability it assigned to the outcome that occurred. Its expected payoff is therefore the negative entropy of its belief, and the expected improvement in that payoff from an observation is exactly the expected information gain. \citet{bernardo1979} showed on this basis that expected information is expected utility when rewards are proper log scoring rules. In the bandit setting, information-directed sampling (IDS) \citep{russo2014} explicitly trades off instantaneous regret against information gain by minimizing the information ratio $\Gamma_t = \delta_t^2 / g_t$, where $\delta_t$ is the expected regret and $g_t$ is the information gain. The key structural difference from EFE is that IDS minimizes the ratio at each step (a relative weighting that adapts to the current belief), while EFE fixes the weight at $w{=}1$ (an absolute weighting in information units under a fixed reward convention). This paper includes an observe-then-commit IDS baseline (Section~\ref{sec:methodology}, with full results in Appendix~\ref{app:ids}).

Of these classical lines, sequential Bayesian experimental design is the closest structural analogue to the problem class we study, since our observe-then-commit structure parallels its selection of experiments (observations) to maximize information about an unknown state before a terminal decision. The $\rho$-POMDP framework operationalizes this connection by setting the belief-dependent term of each observation action to the expected information gain that action yields. Section~\ref{sec:methodology} defines that quantity formally and writes it $I_a(b)$, the expected information gain of observation action $a$ at belief $b$, so that the operationalization reads $\rho = I_a(b)$. The $I(b,a)$ of Figure~\ref{fig:hero} and the bare $I$ of $R + wI$ above denote the same quantity, with the action written in a different place or left implicit. \citet{walraven2024} approach the same connection from the active-inference side, connecting active inference to information gathering in POMDPs and situating EFE-style epistemic value alongside classical POMDP information-seeking objectives.

\paragraph{Active inference and EFE.}
The Active Inference (AIF) framework \citep{friston2010} casts perception and action as variational inference under the free-energy principle. \citet{parr2022} provide a comprehensive textbook treatment. \citet{friston2015} formalized the decomposition of policy value into extrinsic (goal-seeking, the component the Introduction called pragmatic) and epistemic (information-seeking) components, showing that curiosity-driven exploration arises automatically from expected free energy minimization. \citet{dacosta2020} synthesized discrete-state AIF from first principles, showing that the posterior over policies takes the form $Q(\pi) \propto \exp(-\mathcal{G}(\pi))$ where $\mathcal{G}$ is the Expected Free Energy. \citet{parr2019} showed that EFE decomposes into pragmatic value (divergence from preferred observations) and epistemic value (expected information gain), with both arising from a single variational objective and requiring no tunable exploration weight. The same paper compared EFE against a second functional, the generalised free energy, which differs in construction. Where EFE averages a free energy over the observations a policy is predicted to produce, the generalised free energy treats those future observations as unobserved variables to be inferred alongside hidden states, within one free-energy functional spanning past and future time points. Despite different constructions, the two functionals produced closely matching policy posteriors in that paper's simulations. They are not identical in general, and the coefficient result of Proposition~\ref{prop:equivalence} is established here for EFE only.

The mathematical foundations of EFE have been critically examined. \citet{millidge2021} showed that naively extending the variational free energy, the perception-time objective that upper-bounds the surprise of observed data, into the future does not yield exploratory behavior, proposing the Free Energy of the Expected Future (FEEF) as an alternative with clearer mathematical grounding. The literature also contains four algebraically distinct expressions that are all called expected free energy. \citet{champion2024} addressed the ``unification problem'', whether the four EFE formulations in the literature follow from a single root definition, by studying two candidate roots. The two roots trade coverage against justification. All four formulations follow from the first root, but at a price. Under that root only a restricted class of prior preferences over observations is compatible with the generative model's likelihood mapping. The second root can be justified directly as a definition in its own right, but it recovers only two of the four formulations. \citet{devries2025} recast EFE-based planning as variational inference on a generative model augmented with preference and epistemic priors, an alternative derivation intended to enable message-passing implementations that the paper itself leaves to future work. Our own results depend on one of these formulations only. The equivalence of Section~\ref{sec:methodology} is established for the discrete-state EFE whose epistemic term is the expected information gain about the hidden state, evaluated recursively over belief trajectories (the sophisticated-inference form described next). Whether the same coefficient survives under FEEF, the generalised free energy, or the recasting just described is not assumed. Section~\ref{sec:methodology} lists those alternatives, together with one further derivation not discussed in this section, as cases the result leaves open.

Two contemporaneous 2024--2026 papers examine EFE's relationship to other decision-theoretic targets rather than its internal formulation. \citet{wei2024voi} analyze EFE's optimality gap against a reward-only Bayes-optimal policy within a belief-MDP casting, treating the epistemic term as an information-value correction rather than an exact identity with a specific target. \citet{sweeney2026equivalences} survey formal correspondences between EFE minimization and six decision-theoretic frameworks (Bayesian decision theory, resource rationality, optimal control, reinforcement learning, rate-distortion theory, and maximum-entropy inference), a broader unification effort distinct from our narrower exact identity with $\rho$-POMDPs under log scoring at $w{=}1$ (Proposition~\ref{prop:equivalence}).

\paragraph{Sophisticated inference.}
Standard AIF evaluates policies open-loop. It scores whole action sequences under the current belief without conditioning later actions on intermediate, counterfactual observations. \citet{friston2021} introduced sophisticated inference, a closed-loop recursive extension implementing deep tree search over belief trajectories rather than states. Sophistication (maintaining beliefs about future beliefs) enables counterfactual reasoning about the downstream epistemic consequences of actions. \citet{dacosta2020b} proved an optimality result for this recursive scheme in a specific regime. In their setting the preferences over outcomes encode a reward scaled by an inverse temperature $\beta$, called a precision in active-inference vocabulary, so that larger $\beta$ makes the preferences encode the reward more sharply. As $\beta \to \infty$, the zero-temperature limit of precise reward-encoding preferences, the reward term dominates, so that the epistemic terms only break ties among reward-maximizing action sequences. In that limit the recursive scheme recovers a policy that is Bellman-optimal for the reward, for any finite horizon. The result is stated for POMDPs and assumes exact state estimation. Standard AIF, by contrast, recovers such reward-Bellman-optimal policies only for single-step planning. Our recursive EFE agent (Equation~\ref{eq:efe_recursive}) is derived from this framework, adapted to the $\rho$-POMDP commit-action structure. The recursive counterfactual reasoning over future beliefs is preserved. What changes is that our observe-then-commit setting does not involve state transitions between time steps, simplifying the belief-trajectory computation.

\paragraph{Scaling active inference.}
Discrete-state AIF with full policy enumeration is limited to small state--action spaces. Several lines of work address scaling. \citet{fountas2020} combined deep generative models with Monte Carlo tree search for EFE-optimal planning in continuous state spaces. \citet{tschantz2020} developed an RL-compatible objective (the FEEF objective introduced above) that inherits AIF's exploration--exploitation balance while scaling to standard RL benchmarks. \citet{maisto2025} combined active inference with Monte-Carlo tree search for large POMDPs, reproducing state-of-the-art POMDP solutions on RockSample \citep{smith2004}. Our MCTS-EFE variant (Section~\ref{sec:discussion}, Appendix~\ref{app:pomcp}) follows this direction, using EFE as the in-tree observe reward and leaf heuristic within MCTS to extend planning horizons beyond exact tree search.

\paragraph{Exploration, control as inference, and intrinsic motivation.}
The control-as-inference perspective \citep{todorov2007, levine2018} casts maximum-entropy reinforcement learning as probabilistic inference. The correspondence is exact under deterministic dynamics and only variational under stochastic dynamics. The reason is that conditioning on optimality under stochastic transitions would credit the agent with control over chance outcomes, and the variational treatment removes that by holding the transition model fixed. Soft actor-critic \citep{haarnoja2018} and KL-minimizing stochastic optimal control \citep{rawlik2012} are algorithmic instances. \citet{millidge2020} provided a formal comparison of AIF and control-as-inference, tracing their differences to how value is encoded in the generative model. \citet{sajid2021} showed that intrinsic motivation behaviors arise under EFE. Intrinsic motivation methods, including curiosity \citep{schmidhuber1991, pathak2017}, count-based exploration \citep{bellemare2016}, VIME \citep{houthooft2016}, and random network distillation \citep{burda2019}, all require tunable bonus weights. \citet{oudeyer2007} survey and classify the computational forms these intrinsic signals take. Bayes-Adaptive MDPs \citep{duff2002, guez2013} and Bayesian RL more broadly \citep{ghavamzadeh2015} address model uncertainty (unknown dynamics), distinct from our focus on state uncertainty under a known model. Our $\rho$-POMDP formalism makes the connection between EFE and these lines of work precise while fixing a derived information-unit weight under a shared reward convention.

\paragraph{Resource-constrained information acquisition.}
A separate literature treats information acquisition itself as resource-constrained, without reference to $\rho$-POMDPs. Rational inattention \citep{sims2003,matejka2015} takes an information-processing constraint as primitive, imposed as a modeling assumption rather than derived. In \citet{sims2003} the constraint is a Shannon-capacity cap on how many bits per period the decision-maker's chosen action may carry about the state, and in \citet{matejka2015} it is instead a cost charged per unit of information acquired. Constrained MDPs dualize an expected-cost constraint via a Lagrange multiplier \citep{altman1999}, and constrained POMDPs enforce the cumulative-cost bound directly by point-based dynamic programming over belief and admissible-cost pairs \citep{kim2011cpomdp}. Soft Actor-Critic's automatic entropy-temperature tuning \citep{haarnoja2018sac} drives a resource (policy entropy) to a target by dual gradient descent. Sequential Bayesian experimental design \citep{lindley1956,foster2021dad}, met above as the closest structural analogue to our setting, adaptively selects the next experiment to maximize information about a latent parameter. Section~\ref{sec:budget_positioning} positions our budgeted $\rho$-POMDP formulation in detail against each of these four lines, namely rational inattention, constrained MDPs and POMDPs, automatic temperature tuning, and sequential experimental design. None of these approaches gives a designer an operational sensing budget specifically for $\rho$-POMDP planning, exact scale equivariance of the resulting policy for the receding-horizon planner, or a crossing bracket and endpoint mixture at usage gaps. The next section sets up the observe-then-commit $\rho$-POMDP class and the EFE objective that the rest of the paper, including those comparisons, builds on.
